\section{引言}
\label{sec:introduction}

\subsection{研究背景与关键问题}

随着大语言模型在自然语言理解、任务规划和工具调用方面能力提升，金融信息服务正由关键词检索和单轮问答逐步转向能够执行连续研究任务的智能体系统。上市公司研究天然涉及公告、研报、财务报表、股权关系和事件信息等多类数据，并要求在连续追问中保持主体、时间和分析口径一致，因此是 Agentic AI 的典型应用场景。

然而，通用大模型或简单检索增强生成仍难以直接满足真实研究需求：长周期对话容易出现主体与期间漂移；文本、数值和关系数据难以通过单一检索链路协同；财务计算、时间过滤和关系搜索若完全交由模型自由推理，容易产生不可复现错误；同时，金融研究还具有严格的信息披露时点约束，若混用未来数据，会形成事实上的“时间穿越”。因此，金融智能体不仅要能够回答问题，还需要保证研究过程可控、结论有据且能够审计。

\subsection{FinTrace 总体解题思路}

系统采用单一主 Agent 与专业确定性工具协同的架构：Agent 负责请求理解、复杂任务的下一步调查决策和最终解释；工具负责文档检索、财务计算、关系查询和事件整理；Evidence 体系负责将工具结果转化为统一、可追溯的事实对象。

在线执行不采用一次性生成完整工具计划，也不采用无限 ReAct。系统先恢复会话并解析实体、时间和任务，再进行 Pre-Answerability 判断；简单且无歧义的事实请求走确定性 Direct Fast Path，复杂研究任务进入 Evidence-guided Bounded Investigation，根据当前证据缺口逐步决定下一项工具动作。每个动作均需经过 Schema 与能力边界校验，工具执行后再审查 Evidence 是否充分，并据此继续调查、部分回答或停止。

\subsection{主要技术贡献}

本文围绕上市公司研究中的可信执行形成四项主要工作：

\begin{itemize}
    \item 构建以 Evidence Ledger 为核心的证据驱动研究闭环，使工具结果、证据与最终 Claim 之间形成显式映射；
    \item 设计包含可回答性 Gate、确定性快速路径和有界调查循环的金融 Agent Harness，提高工具调用的稳定性与可审计性；
    \item 建立统一主体与时间语义下的多源金融数据体系，区分业务观察时点与信息可得截止时点；
    \item 将文档检索、财务分析、股权分析和事件时间线工具化，并支持 Agent 根据研究目标进行组合调用。
\end{itemize}
